\newchapter{Bebras India \BrochureYear Organizing Committee}

Vipul Shah, Tata Consultancy Services, ACM India / CSpathshala: Lead, Task developemnt, selection, Localization, Challenge Organization, 
Brochure creation

Madhavan Mukund, Chennai Mathematical Institute: Task developemnt, selection

Sonia Garcha, CSpathshala: Task Translations, Challenge Organization, Teacher development and Support, Brochure publication

\newchapter{Acknowledgements}

We would like to thank the international Bebras community for allowing us to use the tasks they have developed. Bebras is a collective effort of many countries.

We would like to thank ACM India's CSpathshala initiative volunteers for localizing and reviewing the tasks and reaching out to schools.

Translation:  Akankshya Das, Akansha Mehrotra, Amit Saha, Amollh Mahajan, 
Anita Kelkar, Ankit Agrawal, Archana Sahu, Arshakh Muhammed, Aryan Gupta, Aveek Chakraborty, 
Bhadresh Udani, Daud Shaikh,
Harshal Bidkar, Hiren Mehta, Isha Choudhary, Kannan Thamarai, Krishna Elluru, Lalitha Anantha,
Manickam Dhayalan, Manivannan Vaithiyanathan, Mihir Bhatt, Mugdha Sathe, Munal Patel, Neeraj Kumar, 
Pachaiyappan Murugan, Payal Dhaigude, Pink Dutta, Prakruti Manal, Prashant Bopte, Priyanka Sen, Punya Muduli,
Rahul Kulkarni, Rahul Pandey, Rajiv Ranjan, 
Sambhu Kumar,  Sandhya Rao,  Saurabh Kelkar, Saranya M, Siddiq K,
Srishti Singh, Subhasis Martha, Ujvala BhumiReddy, Urs Shweta, Venkatesh Alwan, Vipul Shah
Review: Aditi Karandikar, Arun Kadekodi, Krithika Krishnamoorthy, Lavanniya B S, Shruti Shah, Shweta

We would also like to thank the teachers who have put in the effort for their students 
to participate in the challenge. 

The online version of the competition was implemented on cuttle.org. We would like to thank the following for their excellent collaboration:
Eljakim Schrijvers, Justina Oostendorp, Alieke Stijf, Kyra Willekes, Dave Oostendorp: cuttle.org, Netherlands

\vskip 11pt

%\newpage

%\includegraphics[width=0.5\textwidth]{graphics/bebrasindia.pdf}

The 2026 Bebras India competition was organized by ACM India's CSpathshala initiative with support from Tata Consultancy Services. Raspberry Pi Foundation India (https://www.raspberrypi.org/) and 
Educational Initiatives (https://https://ei.study/) provided several tasks from India.

\vskip 11pt

This booklet of tasks was created on \today{} using the \LaTeX{} typesetting system developed by Christian Datzko and the Bebras toolchain developed by Jean Philippe Pellet. 
We extend our immense gratitude to the entire Swiss Bebras team for generously sharing this system with us. 
In particular, we would like to thank Jean Philippe Pellet and Susanne Falkenberg-Thut for their continuous guidance and support in making the system fully operational.

%Note: All links were checked on November 1, \BrochureYear{}.

\vfill

%\begin{wrapfigure}{l}{0pt}\raisebox{-0.46cm}[\height-0.92cm][-0.46cm]{\href{https://creativecommons.org/licenses/by-nc-sa/4.0/}{\includegraphics[width=0.2\textwidth]{graphics/by-nc-sa.pdf}}}\end{wrapfigure}The Bebras tasks are licensed under a Creative Commons Attribution-NonCommercial-ShareAlike 4.0 International License. 

\includegraphics[width=0.2\textwidth]{graphics/by-nc-sa.pdf}

\noindent Copyright \copyright{} 2026 Bebras, Bebras India.\\
This work is licensed under a \href{http://creativecommons.org/licenses/by-nc-sa/4.0/}{Creative Commons Attribution-NonCommercial-ShareAlike 4.0 International License}.


%The authors are listed on page~\pageref{authors}.
